\section{Do gradients over reasoning predict what structural edits do?}
\label{sec:diagnosis}

\paragraph{Design.} Search succeeds or fails for many reasons, so we test the signal directly.
For a task, a model and an incumbent prompt we draw $24$ training questions and let the model
propose label-free rewrites and insertions for every slot ($20$--$21$ edits per pool). Each
predictor is computed on the $24$ questions; the target is the change in greedy accuracy on $200$
held-out questions. Predictors: \greater{}'s objective (Eq.~\ref{eq:greater}), exactly (one pass
per edited prompt) and by one-pass patching; and, as the evaluation-based reference, fresh greedy
accuracy of every edited prompt on $8$ and on all $24$ questions. Eleven pools: a replication on
Llama-3-8B and Qwen3-8B, three tasks and a non-initial prompt state, with protocol and hypotheses
fixed before its data (v3, eight pools), and a prospective test on an untouched task with three
models (TS3). The validity tasks (logical deduction, tracking shuffled objects) are outside the
benchmark tasks of \S\ref{sec:experiments}. Pools are the units of inference: we report how many
pools agree in sign and an exact sign-flip permutation test over pools; question- and
row-resampling bootstrap intervals are bias-corrected and conditional on the pools. The
split-half reliability $\kappa$ of each held-out target bounds the attainable correlation
(roughly by $\sqrt\kappa$); a best-3 regret---the held-out change of the three truly best edits
minus that of a predictor's top three---measures what a shortlist cares about, against the exact
expectation for a uniformly random shortlist.

\begin{table}[t]
\centering
\caption{\textbf{Spearman correlation between each predictor and the held-out accuracy change of
LLM-proposed block edits} (per pool, pooled mean, and pools agreeing in sign with the sign-flip
test). LD/TS: logical deduction (3/7 objects) and tracking shuffled objects (7); TS3: untouched
task (3 objects). L: Llama-3-8B, Q: Qwen3-8B, G: Gemma-2-9B; $^\ast$ non-initial prompt state;
$\sqrt\kappa$: reliability diagnostic of the held-out target.}
\label{tab:validity}
\resizebox{0.95\textwidth}{!}{\begin{tabular}{lrrrrrrrrrr}
\toprule
Predictor & LD7-L & LD7-L$^\ast$ & LD7-Q & LD7-Q$^\ast$ & LD3-L & LD3-Q & TS7-L & TS7-Q & Pooled & Cost (s) \\
\midrule
Answer loss, exact (\greater{}) & -0.55 & -0.48 & -0.25 & +0.24 & -0.41 & -0.37 & -0.66 & +0.03 & -0.31 [-0.31, -0.10] & 70 \\
Answer loss, patching & -0.54 & -0.59 & -0.12 & +0.39 & -0.05 & -0.30 & -0.42 & +0.02 & -0.20 [-0.28, -0.07] & 94 \\
Fork margin, exact & -0.14 & -0.35 & +0.18 & -0.21 & +0.02 & +0.06 & -0.05 & -0.58 & -0.13 [-0.24, +0.04] & 305 \\
Fork margin, patching & +0.17 & -0.23 & +0.11 & -0.31 & -0.19 & +0.02 & -0.13 & -0.52 & -0.14 [-0.24, +0.07] & 308 \\
Fresh accuracy, 8 rows & +0.48 & -0.07 & +0.18 & +0.34 & +0.27 & +0.38 & +0.48 & -0.30 & +0.22 [+0.01, +0.26] & 537 \\
Fresh accuracy, 24 rows & +0.58 & -0.15 & +0.05 & +0.20 & +0.61 & +0.42 & +0.57 & +0.15 & +0.30 [+0.08, +0.33] & 1190 \\
\midrule
Ceiling $\sqrt{\kappa}$ & 0.56 & 0.76 & 0.61 & 0.11 & 0.88 & 0.84 & 0.95 & 0.80 & & \\
\bottomrule
\end{tabular}
}
\end{table}

\paragraph{The fixed-reasoning objective ranks block edits in the wrong direction.}
\greater{}'s objective correlates negatively with the held-out change in $8$ of $11$ pools
(pooled Spearman $-0.26$; sign-flip $p=0.017$); its one-pass estimate tracks the exact values
(mean Spearman $0.79$ between them) and is negative in $9$ of $11$ pools ($-0.20$; $p=0.037$).
Fresh accuracy on the same questions is positive in $9$ ($8$ questions; $+0.28$, $p=0.010$) and
$10$ ($24$ questions; $+0.33$, $p=0.004$) of the $11$ pools. As a shortlist the patched objective
is worse than a random one in $8$ of $11$ pools (best-3 regret $0.091$ vs.\ $0.067$; $p=0.022$);
the exact objective is not distinguishable from random ($0.075$; worse in $7$ of $11$, $p=0.14$),
and on the untouched task its correlation is null ($-0.13$) rather than negative. The direction
is consistent across models and tasks; its size is not. Splitting training questions by whether
the incumbent answers them correctly, and ranking each question's loss changes across edits
before averaging (which removes scale differences between questions), the negative correlation
on wrongly answered questions persists for Llama-3 ($-0.25$ to $-0.61$ in all five pools) but not
for Qwen3 ($+0.25$, $-0.04$, $+0.07$): the unfaithful-read mechanism of
Proposition~\ref{prop:faithful} explains a model-dependent part of the failure, the absence of the
distribution term explains the rest. Two preregistered repairs---a decision margin at reasoning
forks and the objective restricted to correctly answered questions---failed their criteria
(Appendix~\ref{app:prereg}).

\paragraph{Edits rewrite the reasoning, even token edits.} The fixed reasoning is the weak point.
Under a reader that we verified to be deterministic (three re-reads of the same prompt give
token-identical reasoning on $200$ questions; the answer read differs on $0$--$0.5\%$), a single
\greater{}-style token substitution \todo{deterministic re-roll numbers: reasoning changed on X\%,
first divergence after Y tokens, correctness flipped on Z\%}, while net accuracy moves by about
$\pm2$ points; at this granularity no predictor, gradient or evaluation, is distinguishable from
zero on $200$ questions (Appendix~\ref{app:results}).

\begin{table}[t]
\centering
\caption{\textbf{Where an edit's effect lives} (Eq.~\ref{eq:decomp}, estimated on the incumbent's
$16$ samples per training question). Read-off flips: fraction of (candidate, question, sample)
whose hard read of the same reasoning changes; $|A|$, $|D|$: mean absolute read-off and
reasoning-distribution terms per edit (accuracy points$/100$); sd of the log-likelihood ratios.}
\label{tab:readoff}
\small
\begin{tabular}{lrrrr}
\toprule
Pool & Read-off flips & $|A|$ & $|D|$ & sd$(\log w)$ \\
\midrule
LD7, L$^\ast$ & 3.2\% & 0.012 & 0.034 & 43.5 \\
LD7, L & 1.4\% & 0.005 & 0.039 & 19.6 \\
LD3, L & 0.2\% & 0.002 & 0.041 & 19.0 \\
\bottomrule
\end{tabular}

\end{table}

\paragraph{Verification noise.} Self-optimization loops shortlist candidates and verify them by
fresh accuracy on training questions. Simulating one round on each pool with the recorded
per-question correctness, accepting the best of all candidates by fresh accuracy on $8$ questions
\emph{lowers} held-out accuracy by $2.3$ points on average (on $16$ questions $-0.8$, on $24$
$+1.1$), while the best edit of each pool would add $5.5$: on small minibatches verification picks
noise. \greater{} verifies on all its training questions; the structural search of
\S\ref{sec:experiments} does the same.
